% Questions about the manuscript and the talk, and their answers.
% Every number below is read from an artifact in the repository: the dataset
% ../results/egb-extremality.json (the same file manuscript/numbers.tex is
% generated from), ../results/egb-rotating-external-contrasts.json, and the
% cost ledger ../prompts/uso-de-tiempo-y-tokens.md. The code quoted is
% ../src/rotating_bh/{extremality,egb_rotating_bvp}.py.
% Build: tectonic -X compile questions-answers.tex
\documentclass[11pt,a4paper]{article}
\usepackage[T1]{fontenc}
\usepackage[utf8]{inputenc}
\usepackage{lmodern}
\usepackage[margin=2.6cm]{geometry}
\usepackage{amsmath}
\usepackage{booktabs}
\usepackage{array}
\usepackage{xcolor}
\usepackage[colorlinks=true,allcolors=blue!55!black]{hyperref}

\definecolor{accent}{HTML}{1F5FA8}
\newcommand{\q}[1]{\subsection*{\color{accent}#1}}
\newcommand{\ext}{\mathsf{ext}}

\title{Questions on the manuscript and the talk}
\author{}
\date{1 October 2026}

\begin{document}
\sloppy
\maketitle
\vspace{-0.8cm}

\noindent\rule{\linewidth}{0.4pt}
\medskip

\noindent Seven questions, each answered from an artifact in the repository
rather than from recollection. Sections~1 and~2 concern one sentence of
section~4 of the manuscript; sections~3 and~4 concern one sentence of the
speaker notes; section~5 concerns the extremal solutions of
\texttt{arXiv:2303.12471}; section~6 whether our solver could construct
them; and section~7 the mass-curve secants of the result slide. What was \emph{not} checked is stated at the end of each answer
and collected under Provenance.

% ================================================================
\section{What does 356 mean, and are the couplings equispaced?}

\noindent The sentence at issue opens section~4, \emph{Results}:
\begin{quote}
``The dataset contains $356$ accepted solutions at $13$ couplings from
$\alpha=0$ to $\alpha=1/2$ in units $r_H=1$.''
\end{quote}

\q{The 356 is a total over couplings, not a count per coupling}

It is the number of converged black-hole solutions \emph{summed over all
thirteen couplings} --- not $356$ per coupling, and nothing to do with the
number of $\alpha$ values. Each coupling carries one \emph{walk}: a
continuation in the spin at that fixed $\alpha$ (section~\ref{sec:walk-ref}
below). The number $356$ is the total length of the thirteen walks.

\begin{table}[h]
\centering\small
\begin{tabular}{@{}rrlr@{}}
\toprule
$\alpha$ & states & walk & $TJ^{1/3}$ range \\
\midrule
$0$     & $43$ & full   & $1.12\times10^{-5}$ to $8.74\times10^{-2}$\\
$0.005$ & $10$ & narrow & $2.44\times10^{-3}$ to $5.10\times10^{-2}$\\
$0.01$  & $11$ & narrow & $1.21\times10^{-3}$ to $5.14\times10^{-2}$\\
$0.02$  & $40$ & full   & $4.61\times10^{-3}$ to $8.15\times10^{-2}$\\
$0.035$ & $11$ & narrow & $3.76\times10^{-3}$ to $5.13\times10^{-2}$\\
$0.05$  & $45$ & full   & $4.49\times10^{-4}$ to $7.42\times10^{-2}$\\
$0.075$ & $16$ & narrow & $4.38\times10^{-4}$ to $4.75\times10^{-2}$\\
$0.1$   & $45$ & full   & $3.84\times10^{-4}$ to $6.48\times10^{-2}$\\
$0.15$  & $17$ & narrow & $2.57\times10^{-4}$ to $3.90\times10^{-2}$\\
$0.2$   & $45$ & full   & $1.44\times10^{-4}$ to $5.21\times10^{-2}$\\
$0.3$   & $16$ & narrow & $1.06\times10^{-4}$ to $2.59\times10^{-2}$\\
$0.4$   & $13$ & narrow & $5.06\times10^{-5}$ to $1.97\times10^{-2}$\\
$0.5$   & $44$ & full   & $3.34\times10^{-6}$ to $3.34\times10^{-2}$\\
\midrule
\textbf{total} & $\mathbf{356}$ & & \\
\bottomrule
\end{tabular}
\caption{States per coupling, read from \texttt{results/egb-extremality.json}.
The two populations are the two kinds of walk described in section~2.}
\end{table}

The counts are bimodal, and the split is exactly the one the next sentence of
section~4 describes. The six couplings walked over the \emph{whole} spin range,
$\alpha=0$, $0.02$, $0.05$, $0.1$, $0.2$, $0.5$, carry $40$--$45$ states each,
$262$ in total. The seven that start at $q=0.6$ and run only to the
near-extremal end carry $10$--$17$ each, $94$ in total. So
\[
  356 \;=\; \underbrace{262}_{\text{6 full walks}} \;+\; \underbrace{94}_{\text{7 narrow walks}} .
\]
``Accepted'' means the state passed the solver's acceptance gate: the stored
rows carry the diagnostics it is judged on
(\texttt{max\_relative\_tensor\_residual}, \texttt{psi\_window\_spread},
\texttt{mass\_tail\_spread}, and so on). Refused attempts are not in the count.

\q{The couplings are not equispaced}

\begin{center}\small
\begin{tabular}{@{}l@{\;}*{13}{r@{\;}}@{}}
$\alpha$ & $0$ & $.005$ & $.01$ & $.02$ & $.035$ & $.05$ & $.075$ & $.1$ & $.15$ & $.2$ & $.3$ & $.4$ & $.5$\\
gap & & $.005$ & $.005$ & $.01$ & $.015$ & $.015$ & $.025$ & $.025$ & $.05$ & $.05$ & $.1$ & $.1$ & $.1$\\
\end{tabular}
\end{center}

\noindent The spacing is graded, widening by a factor of about $20$ from the
small-$\alpha$ end to the large-$\alpha$ end. That matches what is being
measured: the result is the departure from the first-order value
\[
  M_\ext=\tfrac32\pi^{1/3}J^{2/3}+\pi\,\alpha+\mathcal O(\alpha^2),
\]
so the $\alpha\to0$ end has to be resolved finely enough to see the slope leave
$\pi$, while the large-$\alpha$ end only has to track a slowly varying curve.

The grid is also aligned with its two uses. The six full walks sit at
$0$, $0.02$, $0.05$, $0.1$, $0.2$, $0.5$, and the seven narrow ones interleave
between them. The narrow walks exist only to add resolution to the
extremal-slope curve, which is why they skip the low-spin region entirely.

% ================================================================
\section{What is a ``walk''?}
\label{sec:walk-ref}

A \emph{walk} is one continuation sequence: a chain of solutions at
\emph{fixed} $\alpha$, stepped in the spin, each solve seeded by the previous
one. It is the unit the dataset is organised into --- thirteen couplings,
thirteen walks, $356$ states. What follows is read from
\texttt{src/rotating\_bh/extremality.py}.

\q{Where a walk starts}

\texttt{walk\_to\_extremality(alpha\_gb)} first builds a base solution with
\texttt{predictor\_ladder(start=.60, \dots)}: it climbs in the \emph{coupling}
from the known $\alpha=0$ Myers--Perry solution up to the target $\alpha$ in
steps of \texttt{coupling\_step=.005}, holding the spin at $\Omega_H=0.6$. In
units $r_H=1$ one has $q=r_H\Omega_H=\Omega_H$, which is why the manuscript
calls this $q=0.6$.

\q{Where a walk goes}

From that base it runs \texttt{walk} twice, in both directions in the spin,
with \texttt{spin\_step=.01} in $\Omega_H$:
\begin{itemize}\itemsep1pt
  \item \texttt{direction=-1}, down to $\Omega_H=0$, the static limit;
  \item \texttt{direction=+1}, up until the temperature falls below
        \texttt{floor=1e-6}, the near-extremal end.
\end{itemize}
The walk is the concatenation \texttt{descending + [anchor] + ascending},
sorted in $\Omega_H$. Its length is that coupling's row in the table above.
This is precisely the six-versus-seven split: the couplings given both
directions are the $40$--$45$-state full walks, while the seven narrow ones
were given only the ascending branch from $q=0.6$ --- the region the $T\to0$
extrapolation actually consumes.

\q{Why it is a walk and not a scan}

Every step is Newton-solved from the previous solution as its seed. These
states cannot be solved independently, because near extremality there is no
good initial guess except a neighbour. Three behaviours of the loop follow from
that, each recorded in the source as the answer to a way an earlier version
failed:

\begin{itemize}\itemsep3pt
  \item \textbf{Halve, do not stop.} A Newton failure is treated as a seeding
  accident rather than a physical boundary: the step halves and retries. What
  ends a walk is the step collapsing below \texttt{minimum\_step=1e-5}. The
  docstring puts it plainly --- ``a wall which looked physical was not''.

  \item \textbf{Roll back on a refused state.} A state that misses the
  acceptance gate is not kept as the seed. Otherwise every later solve inherits
  the bad profile and the walk dies far short of where it could reach, which
  the source records as what the first production run did at the two smallest
  couplings.

  \item \textbf{Raise the resolution before shortening the step.} Approaching
  extremality the independent tensor residual grows, and the cure is
  $N=64\to96$, not a smaller step. The resolution that accepted each state is
  stored with it, which is the \texttt{resolution} field in the rows.
\end{itemize}

\q{Why 356 is not simply steps times couplings}

A walk steps far more often than it records. States are kept on a schedule:
whenever the spin has moved by \texttt{spin\_record=.02} \emph{or} the
temperature has changed by a factor of \texttt{temperature\_record=1.3}.
Measuring at every accepted step would spend the run in \texttt{diagnose};
measuring on a fixed spin grid would crowd the samples where nothing changes.
So $356$ counts \emph{recorded, accepted} states, and the number of Newton
solves behind them is substantially larger.

This also explains the temperature coverage in the table of section~1. The
$1.3\times$ temperature trigger makes the samples bunch up automatically as $T$
collapses towards extremality --- which is where \texttt{extremal\_state} needs
them, since it fits only the \texttt{window=12} coldest states.

% ================================================================
\section{What is the ``rotating solver''?}
\label{sec:solver}

\noindent The sentence at issue is in the speaker notes for the cost slide:
\begin{quote}
``Reproducing all known solutions took about four hours; the first new object,
the rotating solver, took five times more \dots''
\end{quote}

\noindent It is the code that finds rotating Einstein--Gauss--Bonnet black
holes: \texttt{src/rotating\_bh/egb\_rotating\_bvp.py}, built in milestone
(\emph{hito}) 4. The ten-minute notes say it at greater length --- ``the first
genuinely new object, the rotating Gauss--Bonnet black hole'' --- which is what
the five-minute line compresses.

\q{What it is}

A spectral collocation solver for the radial boundary-value problem of the
rotating EGB system. From its docstring:
\begin{itemize}\itemsep1pt
  \item Chebyshev--Lobatto nodes; damped Newton with a \emph{complex-step}
  Jacobian and a line search --- the same architecture as the vacuum solver
  \texttt{vacuum\_bvp.py}, deliberately mirrored.
  \item Four compact amplitudes $B,F,H,W$ on a compactified radius, $z=1/r$ and
  $x=1-z$, so that $x=0$ is the horizon $r_H=1$ and $x=1$ is $r\to\infty$.
  \item The horizon row is singular in the generic interior formula --- the
  coefficient matrix of the underlying $[f',b'',h'',w'']$ system degenerates
  there --- so it uses separately derived horizon relations, verified against
  the exact Myers--Perry closed form.
  \item At infinity, $B=1$, $F=1$, $H_x=0$, $W_x=0$.
  \item At $\alpha_{\rm GB}=0$ the interior equations reproduce the vacuum ones
  exactly.
\end{itemize}
Everything in the paper rests on it: the walk of section~\ref{sec:walk-ref}
calls into it at every step, and the $356$ states are its accepted outputs.

\q{Why it was the expensive block}

Hito 4 cost $11.2$ hours and $508$\,M tokens, $31\%$ of the project's time and
$36\%$ of its tokens --- the largest single block. The ``five times'' of the
notes is a token ratio, and it is a comparison across different units:

\begin{center}\small
\begin{tabular}{@{}lrr@{}}
\toprule
Block & Hours & Tokens \\
\midrule
Hitos 0--3, reproducing the known literature & $4.4$ & $101{,}558{,}821$ \\
\textbf{Hito 4, rotating EGB black hole} & $\mathbf{11.2}$ & $\mathbf{507{,}954{,}588}$ \\
\bottomrule
\end{tabular}
\end{center}

\noindent $508/101.6=5.0$. Note that the sentence quotes hours for the first
block and tokens for the second.

The ledger documents why it was costly:
\begin{itemize}\itemsep3pt
  \item \textbf{The direct tensor derivation was symbolically intractable.}
  Christoffel and Riemann for the non-diagonal metric did not terminate. The
  work moved to the reduced action, with equivalence to the full tensor
  verified \emph{numerically} over generic profiles rather than symbolically.

  \item \textbf{Five discarded solver routes} before the one that worked:
  fitting a single horizon parameter; multi-parameter shooting with
  \texttt{least\_squares} and \texttt{solve\_ivp}; \texttt{solve\_bvp} seeded
  from Myers--Perry with an exact analytic Jacobian, which went singular with
  parameters diverging to $10^4$--$10^{22}$; extracting the horizon conditions
  from the compactified form, where \texttt{sp.limit} did not terminate; and
  translating the raw relations into compact variables, which gave correct but
  unrecognisable expressions that failed to reproduce even the vacuum form. The
  route that worked redid the horizon and infinity expansions \emph{directly}
  in the compact variables.

  \item \textbf{Three self-inflicted errors}, each costed: $z$ and $x$ treated
  as independent symbols; the numerator extracted before the series, producing
  a spurious relation; and $f''$, $w''$ omitted from the physical jets, which
  left an order-one tensor residual even on exact Myers--Perry.
\end{itemize}

% ================================================================
\section{Is it a different solver from the one of arXiv:1010.0860, and what
was it used to reproduce there?}
\label{sec:contrast}

\noindent Different in method and independently written. The second half of the
question has the more interesting answer: the rotating solver was never
successfully checked against their \emph{bulk numbers at finite coupling}. The
single comparison that would do that is the one that failed.

\q{The two solvers}

\begin{center}\small
\begin{tabular}{@{}>{\raggedright\arraybackslash}p{.17\linewidth}>{\raggedright\arraybackslash}p{.37\linewidth}>{\raggedright\arraybackslash}p{.37\linewidth}@{}}
\toprule
 & Brihaye, Kleihaus, Kunz, Radu & This project \\
 & \texttt{arXiv:1010.0860} & \texttt{egb\_rotating\_bvp.py} \\
\midrule
Method & \textsc{colsys}, a standard off-the-shelf BVP package (their ref.~[41]: Ascher, Christiansen \& Russell) & Chebyshev--Lobatto spectral collocation, written from scratch \\[3pt]
Nonlinear solve & Newton--Raphson with adaptive mesh selection & Damped Newton, complex-step Jacobian, line search \\[3pt]
Grid & $10^2$--$10^3$ mesh points, adaptive & $N=64$, raised to $96$ on a refused state \\[3pt]
Quoted accuracy & relative $\sim10^{-6}$ & tolerance $10^{-11}$, with per-state residual gates \\[3pt]
Equations & their own reduction & independently rederived via the reduced action, equivalence checked numerically \\
\bottomrule
\end{tabular}
\end{center}

\noindent What the two share is a strategy, not code. They also continued from
the closed-form solutions, ``increasing gradually $\alpha$ or $\Omega_H$'',
which is structurally the walk of section~\ref{sec:walk-ref} together with the
climb in the coupling that seeds it.

\q{What was actually reproduced against them}

The eight external contrasts of
\texttt{results/egb-rotating-external-contrasts.json}, with their sectors,
which are the point:

\begin{center}\small
\begin{tabular}{@{}llc@{}}
\toprule
Check & Sector & Result \\
\midrule
reduced normalisation pinned & convention & pass \\
$\omega_H^2(j^2)$ reduced relation & Einstein-limit & pass \\
each-spin reading of $J$ & convention & pass \\
critical temperature & Einstein-limit & pass \\
Wald entropy matches paper & Gauss--Bonnet & pass \\
$b'(r_H)=2r_H/(r_H^2+\alpha)$ at $\Omega_H=0$ & Gauss--Bonnet & pass, $\sim10^{-14}$ \\
$b',f'$ decrease monotonically in $\Omega_H$ & Gauss--Bonnet & pass \\
\textbf{published ergosurface radii, their \S4.1.1} & Gauss--Bonnet & \textbf{fail} \\
\bottomrule
\end{tabular}
\end{center}

\noindent To these the manuscript adds their near-horizon eqs.~(4.11) and
(4.12), satisfied to $9.17\times10^{-14}$ and $6.66\times10^{-16}$, which is
what pins the coupling conversion $\alpha_{\rm ref.}=4\alpha$.

The structure matters more than the tally. The file's own scope note says the
contrasts were chosen among statements ``that are closed forms or carry more
digits than the section 4.1.1 radii''. Hence:
\begin{itemize}\itemsep3pt
  \item the Einstein-limit agreements \textbf{do not test the Gauss--Bonnet
  terms at all}, and the file labels them apart for that reason;
  \item the two tightest Gauss--Bonnet agreements are against \emph{closed-form}
  statements --- the static horizon-derivative formula at $\Omega_H=0$, which
  carries no rotation, and the near-horizon algebra, which is computed by
  \texttt{near\_horizon.py} and \textbf{does not use the bulk rotating solver};
  \item the only check that compares our rotating bulk solutions at finite
  $\alpha$ against values they obtained \emph{numerically} is the ergosurface
  radius, which gives $1.102101$ against their $1.104$ at $\alpha=1/2$ in our
  convention. The manuscript states that the discrepancy is unresolved and
  declines to offer it as evidence of an erratum.
\end{itemize}

So, stated strictly: the rotating solver reproduces Myers--Perry exactly at
$\alpha=0$, reproduces closed-form EGB statements, and reproduces the
near-horizon sector --- but its finite-$\alpha$ rotating \emph{bulk} output has
one published cross-check, and that one disagrees in the fourth digit.

The manuscript's containment argument is that the ergosurface radius tests the
reconstruction of the metric functions \emph{away from the horizon}, whereas
the result depends on the asymptotic charge extraction, the response solve and
the zero-temperature extrapolation, none of which uses $r_e$. That is a
reasonable argument and it is made in the discussion rather than buried;
whether it is strong enough is a judgement a referee may well press on.

% ================================================================
\section{Why do the extremal solutions of arXiv:2303.12471 not give
$M_\ext$, and with it an independent check?}
\label{sec:kkr}

\noindent In principle they do, and the manuscript says so: scale invariance
makes the extremal branch one curve in their reduced variables, and that curve
fixes $M_\ext(J,\alpha)$ through $\mu=\mu(0)\,j^{-2/3}$ and $y=4x\mu/(3\pi)$
(\texttt{manuscript/main.tex}, introduction). In practice the paper does not
publish the branch in a form that can be compared at the precision of our
results. Four separate reasons.

\q{1. The near-horizon solution carries no mass}

The exact rotating squashed $\mathrm{AdS}_2\times S^3$ solution (their
eqs.~(3.18)--(3.19)) gives $J$, $A_H$ and $S(\text{extremal})$, not $M$. The mass
is an asymptotic charge and the near-horizon geometry does not fix it; this is
also why \texttt{arXiv:1010.0860} obtained only the extremal entropy.

\q{2. The full extremal solutions give $M$, but no numbers are printed}

The directly constructed extremal solutions do give it,
$M=\tfrac{3\pi^2}{4}(a-c_t)$ with $c_t$ read off the large-$r$ tail
(their eq.~(3.15)). The paper prints no table of $c_t$, $M$ or $J$; the branch
appears only as curves in their Fig.~1, ``resulting from the extrapolation of
around one thousand of data points into the continuum''.

\q{3. The figure does not show the projection we would need}

Every quantity in Fig.~1 is normalised by $M$, and the extremal curve appears
only as two separate projections: $a_H$ against $x=c_0\alpha/M$ in the main
panel, and $a_H$ against $j$ in the inset (likewise for $s$ and $t_H$). There is
no $j$ against $x$, which is what fixes $\mu(y)$. Recovering it means matching
the two projections point by point through $a_H$. As far as can be seen at the
resolution inspected, the inset's extremal curve is flat and not monotonic near
$j\approx0.95$--$1$, so that matching is ill-conditioned there. That is where our
branch lives: $j_\ext=[\mu(0)/\mu]^{3/2}$ with $\mu$ rising by $31\%$ gives $j$
from $1$ down to about $0.67$.

\q{4. Our result is a slope, and slopes read off a figure are the weak point}

The headline numbers ($\pi\to1.07\to1.44$, and the ten checks below $10^{-4}$) are
values of $\partial M_\ext/\partial\alpha|_J$. Differentiating a curve read from
a figure, or differencing separately constructed extremal solutions, amplifies
the error. The paper also states no accuracy for the equal-spin extremal
solutions: the only error estimate it gives ($\sim10^{-3}$) is in its section~4,
for the single-plane solutions. Even a perfect reading would have no error bar to
compare against.

\q{What could still be done}

\begin{itemize}\itemsep3pt
  \item \emph{A coarse check is feasible.} The effects are large --- the slope
  drops by a factor of about three and first order overshoots by $2.3\times$ ---
  so a curve good to the percent level could confirm the sign and the overshoot.
  \item \emph{Getting the curve.} The PDF's Fig.~1 contains no raster images, so
  it is drawn as vector curves; the arXiv source may give the vertices exactly,
  as was done for \texttt{arXiv:1010.0860} (\texttt{short-summary/summary.tex},
  ``Reproducing the published family'').
  \item \emph{The clean option.} Ask the authors for their extremal data
  ($c_t$, $c_w$, $a$ at each $\alpha$). The manuscript's discussion already names
  a direct construction of the extremal branch as what would remove the
  extrapolation in temperature.
\end{itemize}

% ================================================================
\section{Can our solver find extremal solutions, and would constructing them
be an independent check?}
\label{sec:direct}

\noindent Not as it stands; a direct construction would be a worthwhile
extension, though not a fully independent one.

\q{Why the present solver cannot reach $T=0$}

The ansatz of \texttt{src/rotating\_bh/egb\_rotating\_bvp.py} is built around a
non-degenerate horizon: $b=(1-z^2)B$ has a simple zero at $r_H=1$, the horizon row
uses regularity relations derived for that case
(\texttt{\_egb\_rotating\_horizon\_generated.horizon}), and the inputs are
$(r_H,\omega_H,\alpha)$. An extremal horizon is a double zero; the expansion there
changes form and these relations no longer apply. That is why the project
approaches $T=0$ from above --- the coldest state has
$TJ^{1/3}\approx3\times10^{-6}$ --- and why the result rests on a fit in
temperature.

\texttt{near\_horizon.py} does solve an extremal problem, but only the
near-horizon $\mathrm{AdS}_2\times S^3$ algebra ($v_1,v_2,v_3,k$), not the full
solution out to infinity. It therefore gives $S_\ext$, not $M_\ext$
(section~\ref{sec:kkr}, point~1).

\q{What the extension would involve}

\begin{enumerate}\itemsep3pt
  \item \emph{A solver for the degenerate horizon}, for instance in the gauge of
  \texttt{arXiv:2303.12471} (horizon at $r=0$, $u=r^2+a^2$): a series expansion
  at the horizon whose leading data come from \texttt{near\_horizon.py}, and the
  same asymptotic charge extraction as now.
  \item \emph{Seeds} from the coldest states of the existing walks, which at
  $TJ^{1/3}\sim10^{-6}$ should already be very close.
  \item \emph{The slope at $T=0$ directly}: the same implicit-differentiation
  response solve applied to the extremal family at fixed $J$, giving
  $\partial M_\ext/\partial\alpha|_J$ in one linear solve, with no limit in
  temperature and no first-law argument at $T=0$.
\end{enumerate}

\q{What it would and would not check}

\begin{itemize}\itemsep3pt
  \item \emph{It removes the weakest link.} The $T\to0$ extrapolation is what
  the error bars describe, and they are the spread over $21$ fits, which cannot
  reveal an error shared by all of them. Comparing $\mu(y)$ and the slope at the
  $13$ couplings with a direct construction would test it properly.
  \item \emph{It is not fully independent.} It would reuse the same derived
  field equations and the same mass extraction, so an error in either would
  carry over. Independence from those needs the authors' data
  (section~\ref{sec:kkr}) or a second-order perturbative calculation; the
  manuscript's discussion names both, with the direct construction, as the
  three ways to sharpen the result.
\end{itemize}

\q{The main technical risk}

In higher-derivative theories the corrections near an extremal horizon can scale
as non-integer powers of the distance to the horizon, so the solution is not
smooth there; \texttt{arXiv:2303.07358} (in \texttt{papers/}) concerns this for
Kerr. If it happens here, the horizon series needs those powers and the
Chebyshev convergence degrades. \texttt{arXiv:2303.12471} states that it uses a
power series at $r=0$, which suggests the equal-spin case is regular, but this was
not checked.

% ================================================================
\section{Do the mass-curve secants suffer from the finite-difference problem?}
\label{sec:secants}

\noindent Yes. The open squares of the result slide are secants
$\Delta\mu/\Delta y$ between neighbouring couplings, and they have the usual
finite-difference trade-off; in this dataset it matters at the smallest
couplings. The difference from the textbook case is the source of the noise:
not machine rounding, but the error of each extrapolated extremal mass.

\q{The two errors}

For neighbouring couplings $y_1,y_2$, with $\mu=M_\ext/J^{2/3}$:
\begin{itemize}\itemsep3pt
  \item \emph{Noise.} Each $\mu(y)$ is a $T\to0$ extrapolation with an error
  $\delta\mu$, so the secant carries roughly
  $(\delta\mu_1+\delta\mu_2)/\Delta y$, which grows as $\Delta y$ shrinks.
  \item \emph{Truncation.} A secant is the slope averaged over the interval;
  evaluated at the midpoint it misses about $\mu'''\Delta y^2/24$, which shrinks
  with $\Delta y$.
\end{itemize}
So there is an optimal spacing, and below it the noise term blows up.

\q{What the data show}

From \texttt{results/egb-extremality.json}; the noise column is
$(\delta\mu_1+\delta\mu_2)/\Delta y$ from the stored \texttt{mu\_spread}, and for
the first interval it reproduces the stored \texttt{mass\_only\_spread}
($0.263$).

\begin{center}\small
\begin{tabular}{@{}lllll@{}}
\toprule
interval in $y$ & $\Delta y$ & $\delta\mu$ at the ends & secant noise & secant \\
\midrule
$0$--$0.005$       & $0.0048$ & $1.3\times10^{-7}$, $1.3\times10^{-3}$ & $\mathbf{0.26}$ & $3.07$ \\
$0.005$--$0.010$   & $0.0051$ & $1.3\times10^{-3}$, $7.7\times10^{-4}$ & $\mathbf{0.40}$ & $2.91$ \\
$0.010$--$0.021$   & $0.0109$ & $7.7\times10^{-4}$, $1.6\times10^{-3}$ & $0.22$ & $2.65$ \\
$0.021$--$0.038$   & $0.0177$ & $1.6\times10^{-3}$, $8.4\times10^{-4}$ & $0.14$ & $2.22$ \\
$0.038$--$0.057$   & $0.0185$ & $8.4\times10^{-4}$, $5.7\times10^{-6}$ & $0.05$ & $1.82$ \\
$y>0.057$          & $0.03$--$0.10$ & $10^{-5}$ to $10^{-6}$ & $\le0.002$ & $1.47\ldots1.39$ \\
\bottomrule
\end{tabular}
\end{center}

\begin{itemize}\itemsep3pt
  \item \emph{Small $y$.} The couplings are closely spaced and their
  extrapolated masses are uncertain at the $10^{-3}$ level, so the secants are
  uncertain by $10$--$40\%$: the blow-up. That the squares land near the curve
  there is supported by the sampled values, not by their error estimates.
  \item \emph{Larger $y$.} Masses are good to $10^{-5}$--$10^{-6}$ and the
  spacing is wider, so the noise is negligible; the main error is the interval
  average, which is why the squares sit slightly off the dots around the minimum.
\end{itemize}

\q{A rough optimum}

Balancing the two errors gives
$\Delta y_{\rm opt}\approx(48\,\delta\mu/|\mu'''|)^{1/3}$. Near $y\approx0.03$ the
slope samples give $|\mu'''|$ of a few hundred; with $\delta\mu\approx10^{-3}$
that puts $\Delta y_{\rm opt}$ near $0.05$, about ten times the spacing used
there, and with $\delta\mu\approx10^{-5}$ near $0.01$. At small coupling the
secants are noise-limited, and halving the spacing would make them worse.

\q{Why the main method avoids it}

The thermodynamic points difference nothing in $\alpha$: $\Psi$ comes from one
linear solve with the Newton Jacobian at fixed $\alpha$, so there is no
$\Delta\alpha$ to choose (backup slide A2). Its only error is the extrapolation
of $\Psi$ itself. That is also largest at small coupling --- about $0.09$--$0.16$
at $y=0.005$--$0.04$, against $0.001$--$0.003$ at large $y$, because those
couplings were walked over a narrow spin range and their coldest states are not
very cold --- but it is still smaller than the secant noise there ($0.11$
against $0.26$ in the first interval). This is the manuscript's argument for
measuring the derivative rather than differencing masses.

% ================================================================
\section*{Provenance}

The per-coupling counts, the $\alpha$ grid and the temperature ranges were read
directly from \texttt{results/egb-extremality.json} and reproduce the generated
macros \texttt{\textbackslash Nstates}~$=356$,
\texttt{\textbackslash Ncouplings}~$=13$,
\texttt{\textbackslash Nmapcouplings}~$=6$,
\texttt{\textbackslash Nnarrowcouplings}~$=7$ and
\texttt{\textbackslash Nmaplist} of \texttt{manuscript/numbers.tex}. The
description of the walk is read from
\texttt{src/rotating\_bh/extremality.py}.

The description of the rotating solver is read from the docstring of
\texttt{src/rotating\_bh/egb\_rotating\_bvp.py}; the costs, the five discarded
routes and the three errors from \texttt{prompts/uso-de-tiempo-y-tokens.md},
sections 2 and 3. The method of \texttt{arXiv:1010.0860} is from its section
4.1.1 and its reference~[41], in the copy under \texttt{papers/}. The contrast
table is read from
\texttt{results/egb-rotating-external-contrasts.json}; the ergosurface numbers
are from the discussion section of \texttt{manuscript/main.tex}.
Section~\ref{sec:kkr} reads eq.~(3.15), section~3.2, eqs.~(3.18)--(3.20),
Fig.~1 and the error statement of section~4.1 of \texttt{arXiv:2303.12471}, in
the copy under \texttt{papers/}; the relations $\mu=\mu(0)\,j^{-2/3}$ and
$y=4x\mu/(3\pi)$ and the $31\%$ growth are from \texttt{manuscript/main.tex}.
That near-horizon data cannot fix an asymptotic charge is standard background,
recalled rather than checked here.
Section~\ref{sec:direct} reads the docstring and residual of
\texttt{src/rotating\_bh/egb\_rotating\_bvp.py}, the module list of
\texttt{src/rotating\_bh/}, the coldest temperature and the $21$ fits from the
speaker notes, and the closing paragraph of the discussion of
\texttt{manuscript/main.tex}. The non-smoothness of extremal horizons in
higher-derivative theories is recalled, not read from
\texttt{arXiv:2303.07358}.
Section~\ref{sec:secants} reads \texttt{mu}, \texttt{mu\_spread},
\texttt{psi\_uncertainty} and \texttt{y} of each extremal record, and
\texttt{mass\_only} and \texttt{mass\_only\_spread} of
\texttt{route\_comparison}, from \texttt{results/egb-extremality.json}.

\medskip
\noindent\emph{Not checked:} no walk was re-run and no solve was performed; the
$356$ stored rows were not individually re-tested against the acceptance
thresholds, and the external contrasts were read from the stored artifact
rather than recomputed. The ratio $508/101.6=5.0$ of section~\ref{sec:solver}
is arithmetic on the ledger table; the notes do not show the division.
For section~\ref{sec:kkr}, nothing was digitised and no vector extraction of
Fig.~1 was attempted; the shape of the inset curve near $j=1$ was judged by eye
from a 110\,dpi rendering of the page.
Section~\ref{sec:direct} is a proposal: no extremal solve was attempted, the
horizon expansion at a double zero was not derived, and whether the equal-spin
extremal horizon is smooth was not checked.
In section~\ref{sec:secants} the noise column and the optimal spacing are
estimates made here, not results of the project; $|\mu'''|$ is read roughly from
differences of the slope samples, and no secants were recomputed at other
spacings.

\end{document}
